\documentclass[10pt,a4paper]{amsart}
\usepackage[margin=19mm]{geometry}
\usepackage{amsmath,amssymb,mathtools}
\usepackage[scr=rsfs,cal=euler]{mathalfa}
\usepackage{xfrac}
\usepackage[unicode,hidelinks,bookmarks=false]{hyperref}
\newcommand{\ie}{i.e.\ }
\newcommand{\cf}{cf.\ }
\newcommand{\resp}{respectively\ }
\newcommand{\Subsection}{\subsection}
\providecommand{\nobreakdash}{\nobreak\hbox{-}}
\setlength{\emergencystretch}{2em}
\multlinegap0pt
\usepackage{amssymb}
\newtheorem{lemma}{Lemma}
\title{\large{\bf Euler--Poincar\'e Dynamics and Control on Lie Groupoids}}
\author{Ghorbanali Haghighatdoost}
\date{Source: arXiv:2509.20044v2 (CC BY 4.0)}
\begin{document}
\maketitle

\noindent\emph{Source and licence.} \large{\bf Euler--Poincar\'e Dynamics and Control on Lie Groupoids}. Version arXiv:2509.20044v2 (CC BY 4.0), distributed by arXiv under the Creative Commons Attribution 4.0 International licence, http://creativecommons.org/licenses/by/4.0/, as recorded in the arXiv licence metadata for this exact version and shown on the abstract page https://arxiv.org/abs/2509.20044v2. Full licence notice: \texttt{licenses/arxiv-2509.20044-CC-BY-4.0.txt}.\par\smallskip

\noindent\textit{Editorial scope.} Counted unit: the article's lemma lemma (source0.tex lines 242--244). The article's complete original proof of it is delivered with it (source0.tex lines 245--247, one \texttt{proof} environment of 3 lines). No mathematics is written, repaired or re-derived: the statement and the proof are the article's own lines, in the article's own notation, and no retained line is substituted or re-flowed.  No further source-local context is needed: every cross-reference of every retained line resolves inside the two delivered spans.  Nothing was omitted from any retained span, so the package carries no omission record.  The retained lines cite 1 works, whose entries are transcribed verbatim from the article's own bibliography source in the e-print and delivered with short editorial labels recorded in the item's reference mapping.  No retained line contains a figure environment, an image-inclusion command or drawing code, so no image is delivered and the package declares no asset dependency.  Editorial formatting only, in addition: an AMS \texttt{amsart} preamble that restates the article's own declarations and macros used by the retained lines, namely the environments \texttt{lemma} on separate counters, together with the standard packages those lines need.  The retained environments print in this document as Lemma 1 in order of appearance, whereas the article prints the same items with the numbering of its own full document.  The complete native source of the article as distributed in the arXiv e-print for this exact version is bundled unchanged as \texttt{sources/185653/source0.tex}.
\begin{lemma}
The associated isotropy groupoid of a regular Lie groupoid $G \rightrightarrows M $ is a Lie groupoid.
\end{lemma}

\begin{proof}
See \cite{Schmeding}.
\end{proof}

\begin{thebibliography}{99}
\bibitem[Schmed]{Schmeding} A. Schmeding, \textit{The Lie group of vertical bisections of a regular Lie groupoid}, Forum Mathematicum, 32, 479--489, (2019).
\end{thebibliography}
\end{document}
